% ══════════════════════════════════════════════════════════════════════════════
%  EXPERIMENT 1 — EVERYDAY FORCES                     (Lesson 3, T3 Week 1)
%  From Lesson 3 of the A4 Forces Lesson Booklet; answer keys 3.* kept.
% ══════════════════════════════════════════════════════════════════════════════
\experimentdiv{Everyday Forces}{Five workstations \textbullet{} Pushing, gravitational, electrostatic, magnetic and buoyancy forces \textbullet{} Observations and inferences}{3}\label{sec:1}
\lessoncodes{VC2S8U14 \textbullet{} VC2S8I02, I03, I05, I06}

\vocabhead
\begin{ctab}
\begin{tabular}{|SW{4.2cm}|W{11.6cm}|}
  \hline\hd & \thead{Word} & \thead{Definition (copy from the slide)} \\ \hline
  & \scaf{\textbf{Inference}} & \scaf{An explanation of what you observed, using what you already know.} \\ \hline
  & \textbf{Electrostatic force} & \cell{3.v2} \\ \hline
  & \textbf{Magnetic force} & \cell{3.v3} \\ \hline
  & \textbf{Buoyancy force} & \cell{3.v4} \\ \hline
  & \textbf{Generalisation} & \cell{3.v5} \\ \hline
\end{tabular}
\end{ctab}

\begin{lipanel}
\inquiryli{Observe carefully and safely, and tell an observation from an inference.}
\understandli{Recognise five everyday forces and describe the direction each one acts in.}
\sctitle
\begin{itemize}[leftmargin=1.4em, itemsep=0pt, topsep=1pt]
  \item I can follow the safety rules at each workstation.
  \item I can write an inference that names the force and its direction.
  \item I can write a generalisation from several observations.
\end{itemize}
\end{lipanel}

\begin{skillsfocus}
\focus{PC3}{Follow each workstation's safety rule.}
\focus{PC6}{Observe and record what happens each time.}
\focus{CO2}{Draw what you saw, with labels and force arrows.}
\focus{PA5}{Spot the pattern and write it as a generalisation.}
\focus{EV2}{Explain what you saw using the forces you know.}
\end{skillsfocus}

% ─── TECHNIQUE ────────────────────────────────────────────────────────────────
\exppart{Technique spotlight}{Blue}

\begin{technique}{Observation or inference?}{PC6}
An \textbf{observation} is what you see, hear or measure: \emph{``The coin hit the floor before the paper.''} An \textbf{inference} explains it using what you already know: \emph{``Air resistance pushed up on the flat paper more than on the coin, so the paper fell more slowly.''} Record your observation first, then write the inference.
\end{technique}

\begin{technique}{Writing a generalisation}{PA5}
A \textbf{generalisation} is a rule that fits \emph{all} your observations, not just one. Start with ``All \dots'' or ``Whenever \dots''. \emph{``Whenever two north poles are brought together, they push apart.''}
\end{technique}

\begin{workedex}{Observation, then inference}
\textbf{Observation:} a balloon rubbed on a jumper sticks to the wall.

\smallskip
\textbf{Inference:} rubbing gave the balloon an electrostatic charge. The charge attracts the wall, so an electrostatic force pulls the balloon \emph{towards} the wall and holds it there against gravity.
\end{workedex}

\begin{wedo}{Observation, then inference}
A fridge magnet holds a sheet of paper on the fridge door. Write one observation and one inference.
\abox{3.wd1}{4cm}
\end{wedo}

% ─── PLAN AND SAFETY ─────────────────────────────────────────────────────────
\exppart{Plan and safety}{Teal}

\textbf{Aim:} To observe five everyday forces --- pushing, gravitational, electrostatic, magnetic and buoyancy forces --- and describe the direction each one acts in.

\smallskip
There are five workstations, A to E. Your group works at one at a time and moves on when your teacher says so, so you may not visit them in order. Answer a workstation's questions before you leave it.

\needspace{14\baselineskip}
\task Complete the risk assessment before you start. \skr{PC3}
\begin{ctab}
\begin{tabular}{|SW{4.0cm}|W{5.2cm}|W{5.4cm}|}
  \hline\hd & \thead{Risk (what could go wrong)} & \thead{Prevention} & \thead{Responding (if it happens)} \\ \hline
  & \scaf{Germs spread from a shared straw (A)} & \scaf{Use only your own straw; bin it as soon as you finish} & \scaf{Wash your hands; tell the teacher if you used someone else's straw} \\ \hline
  & A ball on the floor makes someone slip (A, E) & & \\ \hline
  & Water spills from the tub (C, E) & & \\ \hline
  & A magnet damages a phone or laptop (D) & & \\ \hline
\end{tabular}
\end{ctab}

\begin{infobox}{Safety at every workstation}
Never throw a ball --- a golf ball can hurt someone. Tell your teacher if you are allergic to latex (rubber). Use scissors and the rod only for the task. Paper and Styrofoam burn easily: keep them away from flames.
\end{infobox}

% ─── RESULTS ─────────────────────────────────────────────────────────────────
\exppart{Results: the five workstations}{Amber}

\subsubsection*{Workstation A --- Pushing forces}
\textbf{Materials:} table tennis ball \textbullet{} drinking straws\par
\wssafety{Use only your own straw, and never blow towards anyone's face. Pick up the ball straight away if it falls on the floor.}
\wsstep{1}Use the table tennis ball and a straw to answer these questions.
\pt{a}What happens if you blow on the ball while it is moving towards you?
\abox{3.A1a}{2.2cm}
\pt{b}What happens if you blow on it while it is moving away from you?
\abox{3.A1b}{2.2cm}
\pt{c}What happens if you blow on it while it is moving across in front of you?
\abox{3.A1c}{2.2cm}
\wsstep{2}Play a game of blowball in groups of four, two in each team. Use a straw each, and set up two goals.
\begin{ctab}
\begin{tabular}{|SW{6.0cm}|Z{4.2cm}|Z{4.2cm}|}
  \hline\hd & \thead{} & \thead{Team 1} & \thead{Team 2} \\ \hline
  & \textbf{Number of goals scored} & \cell{3.A2-r1c2} & \cell{3.A2-r1c3} \\ \hline
\end{tabular}
\end{ctab}

\clearpage
\subsubsection*{Workstation B --- Gravitational forces}
\textbf{Materials:} 20 cent coin (or a metal washer) \textbullet{} scissors and paper\par
\wssafety{Use the scissors only to cut the paper, and carry them with the points down.}
\wsstep{1}Cut out a piece of paper the size of the coin.
\wsstep{2}Hold the paper in one hand and the coin in the other, at the same height. Drop them both at the same time.
\pt{a}Write an inference to explain what happened. Say what forces were acting on the coin and the paper, and in which direction.
\abox{3.B2a}{3.4cm}
\wsstep{3}Roll the paper into a ball and drop it and the coin again.
\pt{b}Write an inference to explain what happened.
\abox{3.B3b}{3.0cm}

\clearpage
\subsubsection*{Workstation C --- Electrostatic forces}
\textbf{Materials:} piece of paper \textbullet{} plastic (Perspex) rod and a woollen cloth \textbullet{} running water from a tap\par
\wssafety{Use the rod only for the task. Wipe up any water you spill straight away.}
\wsstep{1}Tear a piece of paper into small bits. Rub the plastic rod briskly on the woollen cloth.
\pt{a}What happens when you bring the rod near the pieces of paper?
\abox{3.C1a}{2.4cm}
\pt{b}The force here is due to an electrostatic charge on the rod. Does the rod need to touch the paper, or can the electrostatic force act over a distance?
\abox{3.C1b}{2.4cm}
\pt{c}Rub the rod on the cloth again, then hold it near a weak stream of running water. What happens?
\abox{3.C1c}{2.4cm}

\clearpage
\subsubsection*{Workstation D --- Magnetic forces}
\textbf{Materials:} 2 bar magnets with their poles marked\par
\wssafety{Keep the magnets away from laptops, phones and other electronic devices.}
\wsstep{1}Hold one magnet in each hand. Slowly bring the end of one close to the end of the other. Write down what happens, then draw the magnets with an arrow for each force.
\abox{3.D1}{1.8cm}
\drawk{3.D1-draw}{4cm}{}
\wsstep{2}Repeat, but use the other end of one of the magnets. Write down what happens.
\abox{3.D2}{2.0cm}
\wsstep{3}Write a generalisation describing the forces between the two magnets. \skr{PA5}
\abox{3.D3}{2.2cm}

\clearpage
\subsubsection*{Workstation E --- Buoyancy forces}
\textbf{Materials:} bucket or tub of water \textbullet{} balls of different types, e.g.\ golf, table tennis, tennis, rubber, Styrofoam\par
\wssafety{Never throw a ball. Don't carry the tub when it is full, and wipe up spills straight away.}
\wsstep{1}Put a table tennis ball into the water. Push it to the bottom and let it go.
\wsstep{2}Do the same with the other balls. What forces act on a ball in water? In which direction do they act?
\abox{3.E2}{3.0cm}
\pt{a}Write a generalisation about what happens to different balls in water. \skr{PA5}
\abox{3.E2a}{3.0cm}

% ─── ANALYSIS ────────────────────────────────────────────────────────────────
\exppart{Analysis and evaluation}{Coral}

\task Complete the summary of the five forces. \skr{EV2}
\begin{ctab}
\begin{tabular}{|SW{3.6cm}|W{2.6cm}|W{3.2cm}|W{5.2cm}|}
  \hline\hd & \thead{Force} & \thead{Contact or non-contact?} & \thead{Direction it acted} & \thead{Where you saw it} \\ \hline
  & \scaf{Pushing (blowing)} & \scaf{Contact (the air touches the ball)} & \scaf{Away from the straw} & \scaf{Workstation A: the ball moved away when blown} \\ \hline
  & Gravitational & & & \\ \hline
  & Electrostatic & & & \\ \hline
  & Magnetic & & & \\ \hline
  & Buoyancy & & & \\ \hline
\end{tabular}
\end{ctab}

\needspace{12\baselineskip}
\task At Workstation B, the flat paper and the screwed-up paper have the same mass. Use your observations to explain why they fell differently.
\abox{3.t2}{5cm}

\needspace{12\baselineskip}
\task Which workstation was hardest to observe accurately? Suggest one change that would make the observations more reliable.
\abox{3.t3}{5cm}

\begin{yr8box}
\begin{keepext}
\exttask A golf ball sinks and a table tennis ball floats, although both are round. Use the idea of the forces acting on each ball to suggest why. Draw a force diagram for each ball in the water.
\drawk{3.x1}{6cm}{}
\end{keepext}
\end{yr8box}

% ─── SKILL CHECK ─────────────────────────────────────────────────────────────
\exppart{Skill check}{NavyBlue}

\begin{skillcheck}
\scrow{PC3}{I followed each workstation's safety rule.}
\scrow{PC6}{I recorded what I saw before explaining it.}
\scrow{CO2}{My magnet drawing has labels and force arrows.}
\scrow{PA5}{I wrote a generalisation that fits all my observations.}
\scrow{EV2}{I explained each observation with a force I know.}
\end{skillcheck}

\task Which skill will you focus on in Experiment 2? Explain what you will do differently.
\abox{3.sc}{5cm}
